% =====================================================================
% Chapter 9 — Discussion
% =====================================================================
\chapter{Discussion}
\label{ch:discussion}

This chapter interprets the results of both routes, discusses the coverage of
evidence across the five object categories (\secref{sec:disc_coverage}), the
threats to validity (\secref{sec:disc_threats}), and the implications of the
project's central claim --- that the two routes answer different questions and
must be judged accordingly (\secref{sec:disc_implications}).

\section{What the results mean together}
\label{sec:disc_together}

Taken together, the results of \secref{ch:resultsA} and \secref{ch:resultsB}
support the dissertation's thesis: the two routes are complementary answers
that succeed precisely where the other route's assumption fails.

\routeA{} demonstrates that \emph{when overlapping views exist}, sonar
partials can be aligned into a self-consistent measured model: 24/24 views per
object were connected by pose-graph optimisation, median measured-to-surface
residuals of $2.9$--$5.0\mm$ were achieved by the variants, and the
Poisson/BPA pair provides an honest map of where the evidence is thin. The
price is completeness: unobserved regions are empty or explicitly trimmed, and
the model is only as good as the registration.

\routeB{} demonstrates that \emph{when only one view exists}, a learned prior
can still produce a complete object representation: the PCN converged on all
runs (final training CD $1.1\times10^{-3}$--$2.3\times10^{-3}$), the epoch
sheets show category-shaped completions, and a real chair partial was
completed at the correct scale. The price is fidelity: the inferred geometry is
a hypothesis from the training distribution, not a measurement.

The interesting result is the \emph{mesh-conditioned dummy}: training \routeB{}
on \routeA{} meshes transferred the measured geometry into the learned prior.
This is the strongest evidence that the routes are partners rather than rivals
--- the geometric route's best output becomes the learned route's training
ground, closing a loop between measurement and inference.

\section{Coverage across categories}
\label{sec:disc_coverage}

Evidence is not uniform across the five categories, and the dissertation is
explicit about that (\tabref{tab:coverage}).

\begin{table}[H]
  \centering
  \small
  \begin{adjustbox}{max width=\textwidth}
  \begin{tabular}{@{}p{1.6cm}p{6.3cm}p{6.3cm}@{}}
    \toprule
    Object & \routeA{} (geometric fusion) & \routeB{} (PCN completion) \\
    \midrule
    Chair & Rich: segments, fused cloud, Poisson + BPA, turntable, view-aware
    and ray-TSDF variants, hybrid & Checkpoint (150~ep) + real-partial
    inference shown \\
    Tyre & Rich: segments, reconstruction, turntable, view-aware and ray-TSDF
    variants, hybrid & No dedicated PCN result in this workspace \\
    Dummy & Good: segments, mesh comparison, turntable & Rich: mesh-conditioned
    training, epochs 1--10, comparisons \\
    Drum & Numerically complete: fused cloud and meshes; no dedicated render &
    Rich: epochs 1--10, contact sheet, comparisons \\
    Net & Numerically complete: fused cloud and meshes; no dedicated render &
    No checkpoint or completion image \\
    \bottomrule
  \end{tabular}
  \end{adjustbox}
  \caption{Result coverage across categories. The pipelines completed
  numerically for all five objects; rendered evidence is richest for
  chair/tyre/dummy. The table is a transparency statement, not a quality
  ranking.}
  \label{tab:coverage}
\end{table}

Two points follow. First, \emph{numerical completion} (files produced, graphs
connected, losses decreased) and \emph{validated quality} (independent
reference, external ground truth) are different things; this dissertation has
much more of the former than the latter, as \suop{} allows. Second, the
category coverage asymmetry is a consequence of where the segmentation is
reliable: the chair and tyre produce compact, scale-consistent clusters,
whereas the net's mesh-like structure and the drum's large planar barrel make
segmentation harder or the renders less informative. Future work should fill
the drum/net evidence gaps deliberately.

\section{Threats to validity}
\label{sec:disc_threats}

\subsection{Internal validity}
The main internal threats concern the correctness of the pipelines themselves:

\begin{enumerate}
  \item \textbf{Registration correctness.} FPFH can confuse repeated parts
  (chair legs, tyre tread), and RANSAC can accept a symmetric but wrong
  alignment. The pipeline mitigates this with the fitness $\geq 0.25$ edge
  gate, MST initialisation, loop closures and edge pruning, but it cannot
  prove correctness without ground-truth poses. The graph-connectivity and
  fitness statistics show \emph{consistency}, not \emph{correctness}.
  \item \textbf{Segmentation leakage.} If the object segment contains seabed
  or rig points, registration fits them too. The elevation gating and
  scale-consistent clustering in \secref{ch:dataset} reduce this, but residual
  contamination is likely for the sparse 10~m cases (which were excluded from
  fusion precisely for this reason).
  \item \textbf{Evaluation metric choice.} Chamfer distance rewards uniform
  coverage and can hide local error; counting mesh vertices rewards pipeline
  completion, not quality. Both metrics are used here with their limitations
  stated explicitly rather than presented as accuracy.
\end{enumerate}

\subsection{External validity}
The largest external threat is that \suop{} provides \emph{no external
ground truth} for geometry: no CAD models, no logged object poses, and no
independent scan of the physical objects. Consequences:

\begin{enumerate}
  \item \routeA{}'s ``residual'' figures are self-consistency measures. A
  surface could sit perfectly on noisy returns and still be wrong about the
  real object (e.g.\ a consistently biased sound speed shifts all ranges).
  \item \routeB{}'s training references are synthetic (procedural or
  \routeA{}-derived). The learned prior is therefore a prior over
  \emph{our models of the objects}, not over the real objects. Completeness of
  the inference cannot be validated on the real data.
  \item The object was moved by hand between \suop{} cases; any motion blur or
  non-rigid deformation during acquisition becomes a bias that neither route
  can model.
\end{enumerate}

\subsection{Reliability and reproducibility}
All experiments used pinned open-source versions (Open3D 0.19, PyTorch, NumPy,
scikit-learn; see Appendix~\ref{app:environment}) and were run from
reproducible batch scripts with logged output. The datasets
(\texttt{dataset\_*.npz}) are deterministic given the seed, and the
registration and meshing manifests are committed as JSON in the project, so
every table in this dissertation can be reproduced from the scripts and
inputs. The main reliability caveat is compute: the pairwise registration cost
($\sim$15--18 minutes per object) and the 150-epoch chair run ($\sim$4.1~h)
limit how many configurations could be explored.

\section{Implications for practice}
\label{sec:disc_implications}

Three practical recommendations follow from the results:

\begin{enumerate}
  \item \textbf{Choose the route by data availability, not fashion.} For
  inspection-grade geometry where a measured answer is required, invest in
  multi-view acquisitions and use \routeA{}. For single-pass deployments where
  a complete model is needed, train a category-specific \routeB{} model.
  \item \textbf{Report uncertainty honestly.} The Poisson/BPA divergence map
  (\routeA{}) and the training-distribution caveat (\routeB{}) should travel
  with the results; an unqualified mesh or completion overstates certainty.
  \item \textbf{Link the routes.} The mesh-conditioned training result shows
  that the geometric route's outputs can bootstrap the learned route's prior.
  This hybridisation is a more productive research direction than treating the
  two routes as competitors.
\end{enumerate}

\section{Summary}
\label{sec:disc_summary}

The results are consistent with the dissertation's central claim: \routeA{}
produces measured, self-consistent but partial geometry when views exist;
\routeB{} produces complete, prior-informed but unverified geometry when only
one view exists; and the two can be chained. The principal limitation --- the
absence of external geometric ground truth in \suop{} --- motivates the
concluding chapter's proposed experiment.
